\section{All determinant classes for Gaussian fourfolds}
\label{sec:gaussian-all-determinants}

\paragraph{HD-S11 (an explicit effective descent from split sixfolds).}
Let \(A\) be a polarized abelian fourfold with integral
\(\mathbb Z[i]\)-action of signature \((2,2)\), whose Rosati
involution induces complex conjugation. Let
\[
 \Delta=\frac1{4!}\int_A\theta_A^4.
\]
Every ordinary Hermitian determinant class is allowed.
There is a reduced irreducible effective surface \(Z\subset A\)
with nonzero Weil projection and
\[
 \boxed{\deg_{\theta_A}Z\le63700992\,\Delta^{12}.}             \tag{S11.1}
\]
For \(\delta t\) integral, \(t\ne0\) in the Weil plane, and
\(M_T=\lceil\sqrt{\int_A t^2}\rceil\), the explicit integral
division bound is
\[
 \boxed{c_t\le61440\,\delta M_T
       \bigl(63700992\,\Delta^{12}\bigr)^2.}                 \tag{S11.2}
\]
The degree-48 and degree-192 bounds for the explicitly specified
principal Gaussian lattices remain sharper on those components.

\paragraph{Proof: an actual split sixfold with the required factor.}
Write \(\Lambda=H_1(A,\mathbb Z)\) and let \(E_A\) be the Riemann
form of \(\theta_A\). The second-variable-linear Hermitian form
\[
 h_A(x,y)=E_A(x,iy)+iE_A(x,y)
\]
has integral Gaussian entries. The identity Hermitian matrix
induces a unimodular alternating form on the underlying basis
\((v_j,iv_j)\). Replacing it by its matrix \(H_A\) multiplies the
alternating determinant by
\(N_{K/\mathbb Q}(\det H_A)=(\det H_A)^2\).
Since the alternating determinant is \(\Delta^2\) and the
Hermitian signature is \((2,2)\), one has
\[
 \det H_A=\Delta.                                         \tag{S11.3}
\]
The freeness of \(\Lambda\) over \(\mathbb Z[i]\) follows from
torsion-freeness and the PID property, so this is an integral
matrix calculation.

Take \(E_i=\mathbb C/\mathbb Z[i]\), with the degree-one
polarization \(\eta=[0]\), and form the actual polarized surface
\[
 B=E_i^2,\qquad
 i\cdot(x,y)=(ix,-iy),\qquad
 \theta_B=\eta_1+\Delta\eta_2.
\]
Its compatible Hermitian matrix is
\(\operatorname{diag}(1,-\Delta)\): the standard integral
alternating form satisfies \(E(1,i)=1\), whereas the conjugate
action replaces \(i\) by \(-i\) in the second factor.
The actual product
\[
 X=A\times B,\qquad
 \Theta=p_A^*\theta_A+p_B^*\theta_B
\]
therefore has signature \((3,3)\), ordinary determinant
\(-\Delta^2\), and polarization product \(\Delta^2\).
The determinant class is \([-1]\), since \(\Delta^2=N(\Delta)\).
HD-S03a supplies an explicit rational splitting into \(U^3\)
by successive isotropic-vector changes of basis. Hence HD-S10a
applies to this actual sixfold. Its boundedness was extended to
the entire fixed lattice moduli component before restriction to
a family, and includes this decomposable fibre.
It gives an effective integral detecting three-cycle \(Y\) with
\[
 \int_Y\Theta^3\le
 D_{6,i}(\Delta^2)=(2\Delta^2)^6\,746496.                   \tag{S11.4}
\]

\paragraph{Proof: four actual divisor maps.}
Use the homomorphisms
\[
 f_1(x,y)=x,\quad f_2(x,y)=y,\quad
 f_3(x,y)=x+y,\quad f_4(x,y)=x+iy
\]
and divisor classes \(L_j=f_j^*\eta\), \(D_j=2L_j\).
The degree-two elliptic line bundle is globally generated,
by Riemann--Roch: it has two sections and a one-dimensional
space after imposing any point. Thus each \(D_j\) is globally
generated.

Relative to the product principal polarization on \(E_i^2\),
the four \(L_j\) have matrices
\[
 \begin{pmatrix}1&0\\0&0\end{pmatrix},\quad
 \begin{pmatrix}0&0\\0&1\end{pmatrix},\quad
 \begin{pmatrix}1&1\\1&1\end{pmatrix},\quad
 \begin{pmatrix}1&i\\-i&1\end{pmatrix}.
\]
The cup pairing, computed as the norm of the determinant of
each pair of row homomorphisms, is
\[
 \left(\int_B L_jL_k\right)=
 \begin{pmatrix}
 0&1&1&1\\
 1&0&1&1\\
 1&1&0&2\\
 1&1&2&0
 \end{pmatrix},\qquad \det=-4.                            \tag{S11.5}
\]
For the \(D_j\) the determinant is \(-1024\).
The four classes therefore form a basis of \(H^{1,1}(B)\),
which has complex dimension four, and their intersection
functionals separate its rational subspace.

For these same effective divisor classes the theorem of the
square gives the nef inequalities
\[
 0\le D_j\le4\theta_B.                                    \tag{S11.6}
\]
For \(j=3,4\), use
\((x+uy)^*\eta+(x-uy)^*\eta=2\eta_1+2\eta_2\),
where \(u=1,i\); the two summands are nef and \(u\) preserves
\(\eta\). The coordinate cases use only \(\Delta\ge1\).

\paragraph{Proof: nonzero contraction and the effective degree.}
The exterior-product map places the Weil plane of \(X\) in
the \((4,2)\) Kunneth component, with complex lines
\[
 \left(\bigwedge^4 V_{A,\sigma}\otimes
       \bigwedge^2 V_{B,\sigma}\right)
 \oplus
 \left(\bigwedge^4 V_{A,\bar\sigma}\otimes
       \bigwedge^2 V_{B,\bar\sigma}\right).
\]
Thus
\[
 T_Y=(e_{W_A}\otimes1)[Y]_{4,2}\ne0.
\]
The projector is a rational Hodge morphism. Since \(W_A\) has
type \((2,2)\) and \([Y]\) has type \((3,3)\), one obtains
\[
 T_Y\in W_A\otimes\bigl(H^2(B,\mathbb Q)\cap H^{1,1}(B)\bigr).
\]
By (S11.5), at least one contraction
\((1\otimes\int_B(-)D_j)T_Y\) is nonzero.
This use of the complete four-dimensional divisor space prevents
cancellation by additional Kunneth components of the arbitrary
effective detector \(Y\).

Choose a general actual divisor in that globally generated class.
Its pullback meets \(Y\) properly or avoids a component, and
\[
 Z_j=(p_A)_*\bigl(Y\cdot p_B^*D_j\bigr)
\]
is an effective integral surface cycle. The cycle-class projection
formula identifies its Weil projection with the nonzero contraction
for the selected \(j\). In particular at least one contributing image
has dimension two.

Put \(\alpha=p_A^*\theta_A\), \(\beta=p_B^*\theta_B\).
Nef mixed intersections are nonnegative, and (S11.6) gives
\[
 \begin{split}
 \deg_{\theta_A}Z_j
 &\le4\int_Y\alpha^2\beta\\
 &\le\frac43\int_Y(\alpha+\beta)^3\\
 &\le\frac43(2\Delta^2)^6\,746496
   =63700992\,\Delta^{12}.
 \end{split}
\]
The factor three is the coefficient of \(\alpha^2\beta\).
A detecting reduced irreducible component of \(Z_j\) has no
larger degree. This proves (S11.1), and HD-D02 gives (S11.2).
The method agrees with the product-and-divisor contraction in
Schoen1998, Proposition 10, printed pp.332--333; every map,
polarization and degree estimate used here has been supplied
explicitly. \(\square\)

\paragraph{HD-S12 (a fixed Gaussian order of conductor \(f\)).}
For the integral action of
\(\mathcal O_f=\mathbb Z+f\mathbb Z[i]\), the same polarization
product \(\Delta\) gives
\[
 \boxed{\delta_W(A,\theta)\le
        63700992\,f^{46}\Delta^{12},}                       \tag{S12.1}
\]
and
\[
 \boxed{c_t\le61440\,\delta M_T
        (63700992)^2 f^{108}\Delta^{24}.}                  \tag{S12.2}
\]

\paragraph{Proof.}
Set \(\Lambda'=\mathbb Z[i]\Lambda=\Lambda+i\Lambda\).
Then \(f\Lambda'\subset\Lambda\), defining the actual quotient
isogeny \(q:A\to A'\) of degree \(d=[\Lambda':\Lambda]\).
The Riemann form \(E'=fE\) is integral on \(\Lambda'\):
mixed terms use the integral \(fi\)-action, and
\(E(ix,iy)=E(x,y)\). It remains a positive Riemann form of Hodge type $(1,1)$.
The corresponding polarization satisfies
\[
 q^*\theta'=f\theta,\qquad \Delta'=\frac{f^4\Delta}{d}.
\]
Apply HD-S11 on \(A'\) and pull back the detecting surface.
Finite flat pullback preserves integral effectivity and nonzero
Weil projection. Its degree is bounded by
\[
 \frac d{f^2}\,63700992\left(\frac{f^4\Delta}{d}\right)^{12}
 =63700992\,\frac{f^{46}}{d^{11}}\Delta^{12}.
\]
Dropping \(d^{11}\ge1\) and taking a detecting reduced irreducible
component proves (S12.1).

For clarity the integral Gaussian extraction at this conductor
is itself explicit. In the notation of HD-D02 use
\[
 \begin{split}
 P_f={}&4[f]^*+4[fi]^*
       -[f(1+i)]^*-[fi(1+i)]^*,\\
 Z_f={}&[f(2+i)]^*+7[f]^*.
 \end{split}
\]
Every displayed endomorphism lies in the given order. On \(H^4\),
\(P_f=f^4(16e_W)\), while the skew-adjoint rotation is \(f^4\)
times the rotation of HD-D02. Its data are
\[
 R_f=16f^4,\quad m_f=8f^4,\quad
 s_f=576f^8,\quad L_f=32f^4.
\]
The exact Gram extraction therefore gives
\(c_t\le61440\,f^{16}\delta M_T D^2\).
Using (S12.1) contributes \(f^{16+2\cdot46}=f^{108}\),
which proves (S12.2). \(\square\)

\paragraph{A further Eisenstein fourfold consequence.}
The identical proof with \(i\) replaced by \(\omega\) uses
the four divisor classes of HD-G10, whose half-divisor Gram
determinant is \(-3\), and the cubic sixfold supplier
\(D_{6,\omega}(q)=(2q)^6\,30375\).
It proves, for every determinant class,
\[
 \delta_W\le2592000\,\Delta^{12},\qquad
 c_t\le4455\,\delta M_T(2592000\,\Delta^{12})^2.
\]
For the order of conductor \(f\), the respective additional
factors are \(f^{46}\) and \(f^{108}\). The actual replacement
of the Hermitian form is
\(E(x,\omega y)-\bar\omega E(x,y)\), and the appended surface
has action \((\omega,\bar\omega)\) and matrix
\(\operatorname{diag}(1,-\Delta)\). Thus its product is
explicitly the split determinant-\([-1]\) sixfold to which
HD-S10a applies.

